\documentclass[11pt]{article}

\usepackage[margin=1in]{geometry}
\usepackage[utf8]{inputenc}
\usepackage[T1]{fontenc}
\usepackage{lmodern}
\usepackage{microtype}
\usepackage{booktabs}
\usepackage{multirow}
\usepackage{array}
\usepackage{siunitx}
\usepackage{amsmath,amssymb}
\usepackage{graphicx}
\usepackage{xcolor}
\usepackage[numbers,sort&compress]{natbib}
\usepackage{caption}
\usepackage{tikz}
\usepackage[colorlinks=true,allcolors=blue!55!black]{hyperref}

\sisetup{detect-weight=true,detect-family=true,group-minimum-digits=4}
\captionsetup{font=small,labelfont=bf}
\newcommand{\best}[1]{\textbf{#1}}
\newcommand{\skl}{\ensuremath{S}}

\title{\bfseries What Transfers Across Cities in Street-Flood\\
Susceptibility? A Four-City Benchmark of Models,\\
Features, and Presence-Only Corrections}

\author{%
  Fatimah Emad Eldin\\
  \small Mityaf --- Barmajan Al Munawwarah 2026\\
  \small \texttt{Fatimah@trouve.works}
}
\date{\today}

\begin{document}
\maketitle

% ══════════════════════════════════════════════════════════ ABSTRACT
\begin{abstract}
\noindent
A city that floods only rarely has no flood records, and therefore cannot train
a flood model. The standard remedy is to fit a model where records exist and transfer
it, but the literature that would justify this is evaluated in a way that cannot
measure whether transfer works: flood-susceptibility studies overwhelmingly
report ROC-AUC on artificially balanced samples under random splits, a design
that yields AUC above $0.90$ almost regardless of method. We benchmark
\num{14} model families on \num{107152} $\SI{200}{\metre}$ street cells across
four United States cities --- New York, Chicago, San Francisco and Dallas ---
carrying \num{131298} geolocated citizen flood complaints, under spatial
blocking at natural prevalence, with leave-one-city-out as the declared primary
endpoint. Prevalence spans \num{13}-fold across the four cities
(\SI{5.4}{\percent} to \SI{69.1}{\percent}), which makes lift incomparable, so we
report skill, the fraction of attainable headroom captured. Three findings
follow. First, model choice matters less than the literature's model-comparison
focus implies: gradient boosting leads, but the best and worst non-trivial
learners differ by only $0.099$ skill, and the ranking inverts entirely between
two and four held-out cities --- logistic regression ranks first of \num{14} at
$n{=}2$ and twelfth at $n{=}4$. Second, terrain is the part that fails to
transfer. The six DEM-derived hydrological features score \emph{below random}
across cities ($-0.0039$) while scoring $+0.1032$ within them; converting each to
a within-city percentile rank recovers them to $+0.0689$ and improves the full
model by \SI{7.4}{\percent}, exceeding the gain from any model substitution.
Third, presence-only corrections cannot help a ranked product: we estimate that
only \numrange{0.49}{0.65} of flood-prone cells are ever reported, but the
Elkan--Noto correction is a monotonic rescaling and changes our ranking metric by
exactly zero. We release the benchmark and argue that evaluation design, not
architecture, is the binding constraint on deploying these models where they are
most needed.
\end{abstract}

\noindent\textbf{Keywords:} urban pluvial flooding; flood susceptibility;
domain transfer; spatial cross-validation; presence-only learning; 311 data

% ══════════════════════════════════════════════════════════ INTRO
\section{Introduction}

\subsection{Context}
Al Madinah receives a millimetre or more of rain on roughly \SI{2.3}{\percent} of
days --- about eight days a year --- and days heavy enough to leave water
standing on a road, above \SI{5}{\milli\metre}, number two to three. The
municipality maintains \SI{5280}{\kilo\metre} of street network across
\num{91} districts and \SI{545}{\kilo\metre\squared}. It has no flood inventory.
Nobody surveys ponding during a storm, so there is no record of where water
stood, and each rain event is therefore managed as though it were the first.

That combination --- rare rain, real consequences, no records --- is not
peculiar to Al Madinah. It describes most arid and semi-arid cities, and it is
precisely the regime in which a data-driven susceptibility model is most useful
and least trainable.

\subsection{The problem}
The obvious response is to fit a model where labels exist and carry it to the
city that has none. Whether that works is an empirical question, and the
published literature does not answer it, for two reasons.

First, almost all flood-susceptibility studies evaluate on \emph{artificially
balanced} samples --- an equal number of flooded and non-flooded points, a ratio
chosen by the analyst --- under \emph{random} train/test splits. Both choices
inflate scores: balancing removes the difficulty that class imbalance creates,
and random splitting in geographic data places test cells adjacent to training
cells. Reported AUC above $0.90$ is routine and carries little information about
deployment.

Second, the field almost never tests transfer at all. Spatial blocking within one
city answers whether a model generalises across neighbourhoods, holding drainage
design, climate and reporting culture constant --- which are exactly the
variables that differ in a new city.

\subsection{Contributions}
\begin{enumerate}\itemsep2pt
  \item A four-city benchmark of \num{14} model families on \num{107152} cells at
    natural prevalence under spatial blocking, with leave-one-city-out declared
    as the primary endpoint before running (\S\ref{sec:results}).
  \item Evidence that cross-city rankings are unstable at small numbers of held-out
    groups: the ranking of \num{14} models inverts between $n{=}2$ and $n{=}4$
    cities, with the model ranked first at $n{=}2$ falling to twelfth
    (\S\ref{sec:instability}).
  \item An ablation showing that terrain features, the scientific core of this
    literature, contribute \emph{negatively} to cross-city transfer while
    contributing substantially within a city, and that a per-city percentile
    normalisation recovers them and beats every model substitution
    (\S\ref{sec:ablation}).
  \item A quantification of reporting incompleteness (\numrange{0.49}{0.65} of
    flood-prone cells reported) together with a proof-by-construction that the
    standard presence-only correction cannot change a ranking metric
    (\S\ref{sec:pu}).
  \item A negative result for out-of-distribution distance as a predictor of
    transfer quality, with the confound identified (\S\ref{sec:ood}).
  \item An evaluation in the label-poor target city itself, on \num{11} flood
    locations assembled by hand from official and field sources, where an
    untrained four-term hydrological composite proves \emph{indistinguishable}
    from the best of the \num{14} transferred models --- a result we argue
    settles the deployment question on parsimony rather than on accuracy
    (\S\ref{sec:target}).
\end{enumerate}

\subsection{Structure}
\S\ref{sec:related} surveys the literature and tiers it by evaluation rigour.
\S\ref{sec:data} describes the four cities, their label provenance and their
exploratory statistics. \S\ref{sec:method} states features, models,
hyperparameters, splits and metrics. \S\ref{sec:results} reports the benchmark
and the ranking instability; \S\ref{sec:ablation} the feature ablation and the
normalisation result; \S\ref{sec:pu} the presence-only analysis;
\S\ref{sec:ood} the out-of-distribution analysis and the error analysis by city.
\S\ref{sec:target} carries every ranker to the target city and evaluates it on
the eleven confirmed flood locations there. \S\ref{sec:discussion} interprets,
\S\ref{sec:limits} states limitations, and \S\ref{sec:conclusion} concludes.

% ══════════════════════════════════════════════════════════ RELATED
\section{Background and related work}\label{sec:related}

\subsection{Why cross-study comparison fails}
Table~\ref{tab:lit} groups prior work by \emph{evaluation rigour} rather than by
model, because the tier --- not the algorithm --- determines whether two numbers
are commensurable. Within a tier, scores are roughly comparable; across tiers
they are not.

\begin{table}[t]
\centering\footnotesize
\setlength{\tabcolsep}{3.5pt}
\caption{Prior work grouped by evaluation rigour. \textbf{T1} balanced classes
with a random split; \textbf{T2} one rigour control (an out-of-area test,
principled negative sampling, or a presence-only framing); \textbf{T3} natural
prevalence and/or spatial blocking; \textbf{T3+} adds cross-domain transfer as
the primary endpoint. Comparing headline scores across tiers is invalid: each
control removes a source of optimism and lowers the attainable score. Dashes
mark values the open version of the paper does not state.}
\label{tab:lit}
\begin{tabular}{@{}llllll@{}}
\toprule
Tier & Study & Area & Pos:Neg & Split & Headline \\
\midrule
T1 & Vemula et al.\ 2025 & French Riviera & 1:1 & random & XGB AUC 0.985 \\
T1 & Zhang \& Guo 2026 & Yancheng & 3/4-quantile & random & XGB AUC 0.938 \\
T1 & Ajin et al.\ 2025 & Buz\u{a}u & balanced & random & CatBoost AUC 0.972 \\
T1 & Bersabe \& Jun 2025 & Seoul & balanced & random & RF AUC 0.902 \\
T1 & Dey et al.\ 2024 & New Orleans & balanced & random & RF AUC 0.85 \\
T1 & Wang et al.\ 2023 & Dalian & --- & semi-sup. & GAT best of set \\
\midrule
T2 & Seleem et al.\ 2022 & Berlin & balanced & random, out-of-area & RF 0.96 / \best{0.80} \\
T2 & Han et al.\ 2024 & Beijing & --- & --- & RF, SVM overfit \\
T2 & Zhang et al.\ 2025 & Zijiang & FR $+$ OC-SVM & --- & \best{$+$24\,pp} from sampling \\
T2 & Li et al.\ 2022 & Guangzhou & PU learning & --- & PU $>$ ANN \\
T2 & Yuan et al.\ 2023 & Harris Cty & --- & by storm & RF AUC 0.860 \\
\midrule
T3 & Agostini et al.\ 2024 & New York & natural & by period & AUC \best{0.642} \\
T3 & Huang et al.\ 2026 & Houston & regression & leave-location-out & $R^2$ 0.663 \\
T3 & Agonafir et al.\ 2022 & New York & count & --- & buildings $>$ elevation \\
\midrule
\textbf{T3+} & \textbf{This work} & \textbf{4 US cities} & \textbf{natural} &
  \textbf{spatial $+$ LOCO} & \textbf{skill 0.186} \\
\bottomrule
\end{tabular}
\end{table}

Three published results establish that the tier matters more than the model.

\textbf{Equal AUC, different maps.} \citet{vakhshoori2018} produced
susceptibility maps scoring \emph{identical} AUC of $0.77$ that disagreed on
where risk lay; maps from one method differing only in random seed disagreed on
at least \SI{29}{\percent} of pixels, and no equal-AUC pair reached
Cohen's $\kappa > 0.61$ against each other. \citet{lobo2008} give the underlying
reason for presence-only models: AUC depends on the extent from which
pseudo-absences are drawn and is therefore not comparable across studies.

\textbf{AUC survives where the map does not.} \citet{seleem2022} tested Berlin
models outside the training area. Their \SI{30}{\metre} ANN retained
AUC $0.80$ while $\kappa$ fell to $0.02$ --- agreement no better than chance.
This is the closest published analogue to deploying in an unlabelled city.

\textbf{Sampling dominates architecture.} \citet{zhang2025} changed only how
non-flood points were selected and moved accuracy by \num{24} percentage points.
No model substitution anywhere in this literature moves a headline score by more
than about three.

\subsection{Metric practice}
The field reports ROC-AUC, accuracy and $\kappa$, almost always on balanced
samples. Of the studies in Table~\ref{tab:lit} we found exactly one reporting
average precision and none reporting lift over prevalence, despite
\citet{saito2015} and \citet{davis2006} establishing that precision--recall is
the informative curve under imbalance. We adopt PR-based reporting throughout.

\subsection{Presence-only structure}
A cell without a complaint is unreported, not verified dry. This is
positive-unlabelled data, and the parallel to species-distribution modelling is
exact \citep{phillips2009,barbetmassin2012}. \citet{li2022} apply PU learning to
urban flood susceptibility directly. \citet{agostini2024} model under-reporting
in the same New York feed we use, finding reporting rates higher in whiter,
wealthier, owner-occupied tracts, and report AUC $0.642$ --- a value close to
ours and far below the balanced-sample literature, which we read as
corroboration that both are measuring the harder, real problem.

% ══════════════════════════════════════════════════════════ DATA
\section{Data}\label{sec:data}

\subsection{Unit of analysis and label provenance}
The unit is a $\SI{200}{\metre}\times\SI{200}{\metre}$ cell restricted to cells
containing more than \SI{50}{\metre} of road. That restriction is
load-bearing: sampling negatives from arbitrary cells would let a model score
well by answering ``is there a road here?'' rather than ``does water stand
here?''

Labels are geolocated citizen complaints from municipal open-data portals,
selected per city by a descriptor that means water standing on a
\emph{street}. The mapping is explicit per city rather than pattern-matched,
because portals do not agree on taxonomy: New York files the phenomenon as
\emph{Street Flooding (SJ)}; Chicago as \emph{Water On Street Complaint}; San
Francisco files it under \emph{Sewer Issues} and distinguishes street from
property only in a secondary field; Dallas splits it across two municipal
departments. Table~\ref{tab:data} reports the resulting datasets.

\begin{table}[t]
\centering\small
\caption{The four labelled cities and the unlabelled target. Prevalence varies
\num{13}-fold, which is the central methodological fact of this benchmark: it
makes lift incomparable across cities (\S\ref{sec:metrics}). Al Madinah supplies
features only and is never trained or tested on.}
\label{tab:data}
\begin{tabular}{@{}lrrrrr@{}}
\toprule
City & Complaints & Cells & Positive & Prevalence & Blocks \\
\midrule
New York       & \num{21287} & \num{25022} & \num{6860}  & \SI{27.4}{\percent} & 56 \\
Chicago        & \num{80469} & \num{37472} & \num{13404} & \SI{35.8}{\percent} & 79 \\
San Francisco  & \num{25774} & \num{4042}  & \num{2795}  & \SI{69.1}{\percent} & 16 \\
Dallas         & \num{3468}  & \num{40616} & \num{2203}  & \SI{5.4}{\percent}  & 90 \\
\midrule
Total labelled & \num{130998} & \num{107152} & \num{25262} & \SI{23.6}{\percent} & 241 \\
\midrule
Al Madinah     & 0 & \num{8647} & 0 & --- & --- \\
\bottomrule
\end{tabular}
\end{table}

\subsection{Exploratory observations}
Two features of Table~\ref{tab:data} drive the analysis that follows.

\textbf{Prevalence is not a property of flooding.} San Francisco is compact and
intensely reported: \num{25774} complaints over \num{4042} road cells, so
\SI{69.1}{\percent} of cells carry at least one. Dallas is sprawling and sparsely
reported: \num{40616} cells and \num{3468} complaints. The observed rate
therefore reflects reporting density and city form at least as much as physical
susceptibility.

\textbf{Reporting is spatially uneven within cities too.} New York complaints
skew hard by borough --- Queens \num{8169}, Brooklyn \num{6744}, Staten Island
\num{3117}, Bronx \num{1757}, Manhattan only \num{1465} --- a pattern the
literature attributes to combined versus separate sewer design rather than to
Manhattan being dry. This motivates \S\ref{sec:pu}.

% ══════════════════════════════════════════════════════════ METHOD
\section{System and experimental design}\label{sec:method}

\subsection{Features}
Nine numeric features, all computable from global open sources so that every one
is available for Al Madinah unchanged. Six derive from the Copernicus DEM
GLO-30 \citep{esa2022copdem}: \texttt{slope}; \texttt{sink}, the depth of a
closed depression after priority-flood filling \citep{barnes2014};
\texttt{logacc}, $\log(1+\cdot)$ of D8 flow accumulation
\citep{ocallaghan1984}; \texttt{twi}, the topographic wetness index
\citep{beven1979}; \texttt{hand}, height above nearest drainage
\citep{renno2008,nobre2011}; and \texttt{relelev}, elevation minus a
$17\times17$ moving mean. Three derive from OpenStreetMap \citep{osm}:
\texttt{road\_density}; \texttt{major\_frac}, the arterial share of road length;
and \texttt{load\_per\_road}, the sole engineered feature, flow accumulation
divided by road length. There are no images and no reflectance bands: every
feature describes \emph{shape}, which is why the set is computable in a city
with no flood history.

\subsection{Models}\label{sec:models}
Fourteen candidates behind one construction interface, so that the evaluation
code is identical for all. Two things are equalised, because otherwise the
comparison measures preprocessing rather than models. \emph{Class imbalance}:
each candidate receives the equivalent of the incumbent's
\texttt{scale\_pos\_weight} in whatever form its library offers ---
\texttt{scale\_pos\_weight}, \texttt{class\_weight}, or an explicit
\texttt{sample\_weight} at fit time. \emph{Feature scaling}: SVM, MLP,
$k$-NN and logistic regression are wrapped with a standardiser, since
\texttt{logacc} and \texttt{road\_density} differ by orders of magnitude; trees
are left unscaled. Boosted models use \num{350} estimators, depth \num{5},
learning rate $0.06$ and subsample $0.85$. Candidates too costly for
\num{107152} rows are trained on a stratified subsample, preserving prevalence,
and the table records that they were.

\subsection{Splits}
Two protocols. \emph{Within-city}: five-fold \texttt{GroupKFold} where the group
is a $0.05^\circ$ ($\approx\SI{5.5}{\kilo\metre}$) spatial block, so no test cell
borders a training cell \citep{roberts2017,ploton2020}. \emph{Leave-one-city-out}
(LOCO): train on all cities but one, test on the held-out city entire. LOCO is
the declared primary endpoint. Within-city blocking cannot answer the question
this paper asks, because it holds drainage design, climate and reporting culture
constant --- the three things that differ in a new city.

\subsection{Metrics}\label{sec:metrics}
We report PR-AUC and, as the primary quantity, \emph{skill}:
\begin{equation}
  \skl = \frac{\mathrm{PR\text{-}AUC} - \pi}{1 - \pi},
  \label{eq:skill}
\end{equation}
where $\pi$ is prevalence. Skill is $0$ for a random ranker and $1$ for a perfect
one. It is required here because lift, $\mathrm{PR\text{-}AUC}/\pi$, is bounded
above by $1/\pi$: at San Francisco's \SI{69.1}{\percent} the ceiling is
$1.45\times$, at Dallas's \SI{5.4}{\percent} it is $18.4\times$. Averaging lift
across these cities would rank models by which cities happen to have low
prevalence. We report lift alongside for continuity with the literature, and
accuracy only beside its trivial baseline: a classifier predicting ``no flood''
everywhere achieves \SI{72.6}{\percent} accuracy in New York while finding
nothing. Every figure is a mean over five random seeds with standard deviation
reported; deterministic learners are run once and marked.

% ══════════════════════════════════════════════════════════ RESULTS
\section{Results}\label{sec:results}

\subsection{Model comparison}
Table~\ref{tab:board} gives the benchmark. Gradient boosting occupies the top
four places; CatBoost leads at $\skl = 0.1855 \pm 0.0006$ against XGBoost's
$0.1726$, a \SI{7.5}{\percent} relative gain far outside seed variance.

\begin{table}[t]
\centering\small
\caption{Four-city leave-one-city-out benchmark, mean over five seeds
($\pm$ s.d.; deterministic learners run once). Ranked on skill
(Eq.~\ref{eq:skill}). Note the accuracy column: the trivial always-negative
classifier scores \emph{highest} of all. Across the thirteen non-trivial
learners, accuracy and skill are uncorrelated (Spearman $\rho = +0.047$,
$p = 0.88$): accuracy carries essentially no information about ranking
quality here. P@100 is the fraction of the 100 highest-ranked cells that flood.}
\label{tab:board}
\begin{tabular}{@{}lS[table-format=1.4]@{\,}l S[table-format=1.3] S[table-format=1.4] S[table-format=1.3] S[table-format=1.3] S[table-format=1.3]@{}}
\toprule
Model & \multicolumn{2}{c}{Skill (LOCO)} & {Lift} & {Skill (within)} & {ROC} & {Acc.} & {P@100} \\
\midrule
CatBoost            & \best{0.1855} & {\scriptsize$\pm$0.0006} & 1.310 & 0.2326 & 0.649 & 0.585 & \best{0.523} \\
HistGradientBoosting & 0.1768 & {\scriptsize$\pm$0.0027} & 1.295 & 0.2202 & 0.645 & 0.584 & 0.523 \\
XGBoost             & 0.1726 & {\scriptsize$\pm$0.0012} & 1.286 & 0.2198 & 0.643 & 0.587 & 0.506 \\
LightGBM            & 0.1709 & {\scriptsize$\pm$0.0020} & 1.285 & 0.2145 & 0.642 & 0.586 & 0.511 \\
Extra Trees         & 0.1673 & {\scriptsize$\pm$0.0004} & 1.265 & 0.1972 & 0.639 & 0.601 & 0.470 \\
MLP                 & 0.1511 & {\scriptsize$\pm$0.0175} & 1.283 & 0.1799 & 0.638 & 0.641 & 0.457 \\
Random Forest       & 0.1482 & {\scriptsize$\pm$0.0013} & 1.246 & 0.1854 & 0.627 & 0.621 & 0.458 \\
SVM (RBF)           & 0.1467 & {\scriptsize$\pm$0.0038} & 1.268 & 0.1853 & 0.623 & 0.634 & 0.468 \\
$k$-NN              & 0.1422 & & 1.238 & 0.1657 & 0.615 & 0.631 & 0.435 \\
Decision tree       & 0.1406 & & 1.252 & 0.1785 & 0.636 & 0.568 & 0.405 \\
Logistic (L1)       & 0.1354 & & 1.258 & 0.1794 & 0.621 & 0.561 & 0.480 \\
Logistic (L2)       & 0.1353 & & 1.258 & 0.1793 & 0.621 & 0.561 & 0.480 \\
Naive Bayes         & 0.0861 & & 1.176 & 0.1210 & 0.589 & 0.626 & 0.405 \\
\midrule
Always negative     & 0.0000 & & 1.000 & 0.0000 & 0.500 & \best{0.656} & 0.020 \\
\bottomrule
\end{tabular}
\end{table}

Two observations temper that headline. The spread between the best and worst
non-trivial learner is $0.099$ skill, and between the best and the fifth is
$0.018$ --- a narrow band consistent with the \numrange{0.01}{0.03} AUC
separations typical of this literature. And the always-negative row scores the
highest accuracy in the table while attaining P@100 of $0.020$, which is the
clearest possible demonstration that accuracy cannot be used here.

\subsection{Ranking instability at small \texorpdfstring{$n$}{n}}\label{sec:instability}
We first ran this benchmark with two cities. Table~\ref{tab:flip} compares the
rankings. They are not merely noisy; they are close to unrelated.

\begin{table}[t]
\centering\small
\caption{Rank of each model on cross-city skill with two held-out cities versus
four. Computed identically, on the same metric --- the difference is the number
of cities alone. The model ranked first at $n{=}2$ ranks twelfth at $n{=}4$.}
\label{tab:flip}
\begin{tabular}{@{}lcc@{\hskip 2em}lcc@{}}
\toprule
Model & $n{=}2$ & $n{=}4$ & Model & $n{=}2$ & $n{=}4$ \\
\midrule
CatBoost      & 5  & \best{1} & Random Forest & 13 & 7  \\
HistGB        & 6  & 2 & SVM (RBF)     & 3  & 8  \\
XGBoost       & 9  & 3 & $k$-NN        & 12 & 9  \\
LightGBM      & 10 & 4 & Decision tree & 8  & 10 \\
Extra Trees   & 11 & 5 & Logistic (L1) & \best{1}  & 11 \\
MLP           & 4  & 6 & Logistic (L2) & 2  & 12 \\
              &    &   & Naive Bayes   & 7  & 13 \\
\bottomrule
\end{tabular}
\end{table}

We verified that this is not an artefact of changing metric: rankings computed
under lift and under skill agree closely \emph{within} each city count and
diverge \emph{across} them. Two held-out cities produce two transfer numbers, and
their mean is not an estimate of anything stable. We therefore treat any
cross-domain ranking computed on fewer than roughly five held-out groups as
provisional, including our own.

% ══════════════════════════════════════════════════════════ ABLATION
\section{Ablation}\label{sec:ablation}

\subsection{What each feature is worth}
Table~\ref{tab:abl} removes one feature at a time. Gain importance and ablation
disagree sharply, which is the expected behaviour of two quantities that answer
different questions: importance reports what a fitted model \emph{used};
ablation reports what is \emph{lost} without it.

\begin{table}[t]
\centering\small
\caption{Leave-one-feature-out, CatBoost, three seeds, LOCO skill. Baseline
(all nine) is $0.1858$. Gain importance is shown for contrast: \texttt{sink},
\texttt{hand} and \texttt{relelev} together account for \SI{31.6}{\percent} of
gain yet removing any of them \emph{improves} transfer.}
\label{tab:abl}
\begin{tabular}{@{}lS[table-format=1.4]S[table-format=+1.4]S[table-format=2.1]@{}}
\toprule
Removed & {Skill} & {$\Delta$} & {Gain \%} \\
\midrule
\texttt{major\_frac}    & 0.1488 & -0.0370 & 12.9 \\
\texttt{road\_density}  & 0.1797 & -0.0061 & 26.0 \\
\texttt{slope}          & 0.1850 & -0.0008 & 6.2 \\
\texttt{load\_per\_road}& 0.1854 & -0.0004 & 10.2 \\
\texttt{logacc}         & 0.1855 & -0.0003 & 7.1 \\
\texttt{twi}            & 0.1856 & -0.0003 & 6.0 \\
\texttt{sink}           & 0.1862 & +0.0004 & 11.7 \\
\texttt{hand}           & 0.1882 & +0.0024 & 10.3 \\
\texttt{relelev}        & 0.1886 & +0.0028 & 9.6 \\
\bottomrule
\end{tabular}
\end{table}

\subsection{Terrain does not transfer}
Splitting the features into their two natural blocks makes the pattern
unambiguous (Table~\ref{tab:block}). The six DEM-derived hydrological features
--- the scientific core of this literature --- score \emph{below random} across
cities while scoring $+0.1032$ within them. The three OpenStreetMap features
retain \SI{98}{\percent} of their within-city skill when transferred.

\begin{table}[t]
\centering\small
\caption{Feature blocks, evaluated within-city (spatial blocking) and across
cities (LOCO), CatBoost. Terrain carries real signal inside a city and
\emph{negative} signal across cities; urban form transfers almost intact.}
\label{tab:block}
\begin{tabular}{@{}lcS[table-format=1.4]S[table-format=+1.4]@{}}
\toprule
Feature block & $d$ & {Within-city \skl} & {Cross-city \skl} \\
\midrule
All nine                      & 9 & 0.2324 & +0.1858 \\
DEM (terrain) only            & 6 & 0.1032 & -0.0039 \\
OSM (urban form) only         & 3 & 0.1821 & +0.1788 \\
\bottomrule
\end{tabular}
\end{table}

\subsection{Recovering terrain by per-city normalisation}
A terrain feature is an absolute measurement whose meaning is local: two degrees
of slope is flat in San Francisco and steep in Chicago. A decision boundary
learned in one city is therefore not merely weaker in another, it is misplaced.
If that is the mechanism, replacing each terrain value by its percentile rank
\emph{within its own city} should repair it, since ``steeper than
\SI{80}{\percent} of this city'' carries the same meaning everywhere. Crucially,
a percentile rank requires no labels and is therefore computable for
Al Madinah.

\begin{table}[t]
\centering\small
\caption{Per-city percentile normalisation, CatBoost, three seeds, LOCO skill.
Normalising the terrain block converts it from below-random to genuinely
predictive and improves the full model by \SI{7.4}{\percent} --- a larger gain
than any model substitution in Table~\ref{tab:board}. Normalising the urban
features as well \emph{hurts}, consistent with road density in
\si{\kilo\metre\per\kilo\metre\squared} already meaning the same thing in every
city.}
\label{tab:norm}
\begin{tabular}{@{}lS[table-format=+1.4]@{\,}lS[table-format=+1.3]S[table-format=+1.3]S[table-format=+1.3]S[table-format=+1.3]@{}}
\toprule
Configuration & \multicolumn{2}{c}{Skill} & {NYC} & {Chi.} & {SF} & {Dal.} \\
\midrule
DEM only, raw               & -0.0039 & {\scriptsize$\pm$0.0003} & +0.004 & +0.027 & -0.058 & +0.010 \\
DEM only, rank-normalised   & +0.0689 & {\scriptsize$\pm$0.0025} & +0.042 & +0.073 & +0.157 & +0.004 \\
All nine, raw               & +0.1858 & {\scriptsize$\pm$0.0007} & +0.086 & +0.188 & +0.441 & +0.027 \\
All nine, DEM normalised    & \best{+0.1995} & {\scriptsize$\pm$0.0013} & +0.115 & +0.178 & +0.483 & +0.022 \\
All nine, all normalised    & +0.1684 & {\scriptsize$\pm$0.0014} & +0.094 & +0.169 & +0.390 & +0.021 \\
\bottomrule
\end{tabular}
\end{table}

% ══════════════════════════════════════════════════════════ PU
\section{Presence-only analysis}\label{sec:pu}

Our negatives are unverified. Following \citet{elkan2008} we estimate
$c = P(\text{reported} \mid \text{flood-prone})$ from held-out positives under
spatial blocking (Table~\ref{tab:pu}). Roughly half of flood-prone cells are
never reported, which approximately doubles implied true prevalence.

\begin{table}[t]
\centering\small
\caption{Estimated labelling propensity $c$ and implied true prevalence, three
seeds. San Francisco's implied prevalence exceeds unity, which is impossible and
therefore diagnostic: the \emph{selected completely at random} assumption behind
the estimator fails there, consistent with \citet{agostini2024} finding
reporting rates that depend on tract demographics.}
\label{tab:pu}
\begin{tabular}{@{}lS[table-format=1.3]@{\,}lS[table-format=1.3]S[table-format=1.3]@{}}
\toprule
City & \multicolumn{2}{c}{$c$} & {Observed $\pi$} & {Implied true $\pi$} \\
\midrule
New York       & 0.547 & {\scriptsize$\pm$0.0004} & 0.274 & 0.501 \\
Chicago        & 0.552 & {\scriptsize$\pm$0.0001} & 0.358 & 0.648 \\
San Francisco  & 0.649 & {\scriptsize$\pm$0.0010} & 0.691 & \best{1.065} \\
Dallas         & 0.488 & {\scriptsize$\pm$0.0011} & 0.054 & 0.111 \\
\bottomrule
\end{tabular}
\end{table}

The correction changes our results by exactly nothing:
standard $0.185810$, Elkan--Noto $0.185810$, bagging-PU $0.185779$. This is not
an empirical null but a structural one. Elkan--Noto divides every score by a
single constant; a strictly monotonic transform cannot alter any ranking, and
PR-AUC, skill and precision@$K$ are all ranking statistics. \emph{Any}
presence-only correction of this form is incapable of moving these metrics.
Its value is diagnostic: it quantifies how much of the phenomenon is invisible,
and where the estimator itself breaks.

We further tested whether a second, related label --- New York's \num{50821}
catch-basin and manhole-overflow complaints, denoting drainage \emph{failure}
rather than terrain ponding --- is separable with these features. It is not. A
model fitted on ponding scores $0.3187$ on drainage against a same-target
$0.3242$, retaining \SI{98}{\percent}; \SI{81.5}{\percent} of ponding cells are
also drainage cells. The two phenomena are indistinguishable to this feature
set, which is a direct argument that drainage-infrastructure data, absent here,
is the highest-value missing input.

% ══════════════════════════════════════════════════════════ OOD
\section{Error analysis and distributional distance}\label{sec:ood}

\subsection{Where the model fails}
Per-city LOCO skill varies by an order of magnitude: San Francisco $+0.441$,
Chicago $+0.188$, New York $+0.086$, Dallas $+0.027$. The ordering tracks
prevalence almost exactly ($r = +0.977$, $p = 0.023$). Dallas --- sprawling,
sparsely reported, low prevalence --- is both the hardest case and the closest
structural analogue to Al Madinah, and this is the single most consequential
number in the paper for the deployment it motivates.

\subsection{A negative result on OOD distance}
We hypothesised that distributional distance from the training pool would
predict transfer degradation, and could therefore forecast performance in an
unlabelled city. It does not. Distance correlates \emph{positively} with skill
(Mahalanobis $r = +0.881$; PSI $r = +0.993$), meaning the most distant city
transfers best. The relationship is confounded: prevalence predicts skill at
$r = +0.977$ and also correlates with PSI at $r = +0.946$, and with four groups
there are no degrees of freedom to separate them. We report the hypothesis as
refuted rather than reframing it.

What distance can still state is coverage, which is a claim about extrapolation
rather than accuracy. Against all four cities, \SI{5.8}{\percent} of Al Madinah's
\num{8647} cells lie beyond the training 99th percentile of Mahalanobis distance,
and mean PSI is $0.594$ --- a major shift. Per feature, the shift is
concentrated exactly where \S\ref{sec:ablation} showed transfer fails:
\texttt{twi} $1.361$, \texttt{slope} $1.296$, \texttt{road\_density} $0.850$,
against \texttt{hand} $0.053$ and \texttt{major\_frac} $0.007$, which are stable.
The features that do not transfer are the features that shift.

% ══════════════════════════════════════════════════════════ DISCUSSION
\section{Evaluation in the target city}\label{sec:target}

Every result so far is measured on US labels and argued to transfer. This
section measures in the target city itself. It is the weakest evidence in the
paper by an order of magnitude, and we report it in full because a transfer
claim that is never checked against the target is not a claim at all.

\subsection{Assembling ground truth where none is published}

Al~Madinah has no municipal flood register. We built one from resident
recollection with map pins, posts located on X, and archive reports from Saudi
newspapers and the state press agency, corroborating every date against two
independent rainfall reanalyses (ERA5 and NASA~POWER, $\ge 1$\,mm within three
days). This yields 51 dated reports.

Only a fraction carry usable geometry, and the reason is worth recording
because it bounds what any such study can do. Seventeen reports name a date and
nothing more (``the streets of Madinah''); four name only a district. Of the
thirteen that name a street or landmark, four cannot be located even when
OpenStreetMap is searched bilingually --- \emph{Sal\=ah ad-D\={\i}n al-Ayy\=ub\={\i}}
is present only as the Latin string ``Salahedin Al Ayubi Road'', so an
Arabic-only query returns nothing. The binding constraint is that
\textbf{791 of 18{,}320 ways (4\%) in our Madinah extract carry any name tag at
all}.

Eleven distinct locations survive. We stratify them by coordinate provenance,
which turns out to determine the result:
\emph{exact} (1, from our own OSM tunnel geometry),
\emph{precise} (3, a resident placed the pin),
\emph{geocoded} (4, the centroid of a named road, spanning 4.3--16.6\,km),
and \emph{approximate} (3, inferred from post text).

\subsection{Protocol}

The eleven segments are positives; the remaining 30{,}326 are presumed
negatives. Models are the saved fits from Section~\ref{sec:models}, trained on
130{,}998 US complaints and never exposed to a Saudi label. Madinah features are
percentile-transformed against the Madinah \emph{grid} distribution rather than
its road segments, because the models learned percentiles within a gridded city;
scoring segments against segments changes the reference population and cost
$0.02$ AUC in our tests.

Prevalence is $0.036\%$, which disqualifies three familiar metrics. Predicting
``no flood'' everywhere scores $99.96\%$ accuracy. With eleven positives,
precision@100 cannot exceed $0.11$ for any ranker, and $F_1$ inherits that
ceiling. We therefore lead with AUC, which is invariant to the label count, and
with two statistics that are robust to a model's overall calibration: the
\emph{lift} of the mean predicted probability on the eleven over the model's own
city-wide median, and a sign test on how many of the eleven exceed that median.

\subsection{Results}

\begin{table}[t]
\centering
\small
% Emitted by make_target_table.py from the results file, so this table
% cannot drift from the experiment it reports. Re-run the script after
% any change to the results; --check fails if it is stale.
% GENERATED by paper/make_target_table.py -- do not edit by hand.
% Source: research/results_groundtruth_metrics.json
\begin{tabular}{lrrrcr}
\toprule
Ranker & AUC & AP & Lift & Above median & Trustworthy \\
\midrule
Hydrology (untrained) & \textbf{0.685} & \textbf{0.0032} & 1.37 & 8/11 & \textbf{4/4} \\
MLP                   & 0.674 & 0.0007 & \textbf{1.44} & 8/11 & \textbf{4/4} \\
LightGBM              & 0.644 & 0.0007 & 1.06 & 7/11 & 2/4 \\
CatBoost              & 0.642 & 0.0006 & 1.04 & 8/11 & 3/4 \\
HistGradientBoosting  & 0.640 & 0.0010 & 1.04 & 8/11 & 3/4 \\
XGBoost               & 0.635 & 0.0007 & 1.07 & 7/11 & 3/4 \\
Extremely Rand.\ Trees& 0.573 & 0.0009 & 1.06 & 7/11 & 3/4 \\
Random Forest         & 0.537 & 0.0005 & 1.02 & 6/11 & 1/4 \\
$k$-NN                & 0.528 & 0.0006 & 1.05 & 7/11 & 3/4 \\
Decision tree         & 0.515 & 0.0004 & 0.99 & 1/11 & 0/4 \\
Always-negative       & 0.500 & 0.0004 & --- & 0/11 & 0/4 \\
SVM-RBF               & 0.488 & 0.0004 & 0.96 & 6/11 & 2/4 \\
Logistic (L1/L2)      & 0.478 & 0.0014 & 0.96 & 5/11 & 3/4 \\
Naive Bayes           & 0.431 & 0.0005 & 0.72 & 5/11 & 2/4 \\
\bottomrule
\end{tabular}

\caption{Inference on eleven confirmed flood locations in Al~Madinah. Lift is
the mean predicted probability on the eleven divided by the model's own
city-wide median. ``Above median'' counts how many of the eleven exceed that
median (chance is 5.5/11); ``Trustworthy'' is the same restricted to the four
\emph{exact} and \emph{precise} coordinates. Random AUC is $0.500$ and random AP
is $0.00036$.}
\label{tab:target}
\end{table}

Three findings, in decreasing order of confidence.

\paragraph{The untrained hydrological composite matches the best trained model.}
A fixed formula over four terrain indices ---
$\mathrm{pct}(\mathrm{TWI}) + \mathrm{pct}(\mathrm{sink})
- \mathrm{pct}(\mathrm{HAND}) - \mathrm{pct}(\mathrm{slope})$, equal weights,
nothing fitted --- reaches AUC $0.685$ against $0.674$ for the MLP, the strongest
of the fourteen. We deliberately do not call this a win. The gap is $0.011$ at
$n{=}11$; the MLP attains a \emph{higher} lift ($1.44$ against $1.37$); both place
$8/11$ above their own median and all four trustworthy coordinates above it; and
their permutation $p$-values ($0.016$ and $0.023$) are not distinguishable at this
sample size. The defensible statement is that the two are indistinguishable here.

That is the result worth reporting, because the two are not equally expensive.
The MLP requires \num{131298} labelled complaints from four foreign cities, a
training pipeline, and a frozen preprocessing contract; the composite requires
four terrain rasters and a sign convention, can be recomputed by hand, and has no
training distribution to drift away from. Where a supervised transfer merely
\emph{matches} an untrained physical index, parsimony decides --- and a gradient
boosting family that the same benchmark ranks first across US cities
(\S\ref{sec:results}) trails both here, which is itself consistent with
Sections~\ref{sec:ablation} and~\ref{sec:ood}: the learned models lean on road
density, a proxy for \emph{reporting} in complaint data, and Madinah's ground
truth was not assembled from complaints.

\paragraph{Raw hit counts are misleading at this calibration.} Ten of eleven
segments receive $p \ge 0.5$ from CatBoost, which sounds decisive until one
observes that CatBoost assigns the \emph{median} Madinah street $0.734$. The
decision tree makes the point sharply: it returns $p \ge 0.5$ for 11/11 while
having lift $0.99$ and one of eleven above its own median. Any evaluation on a
handful of positives should report lift or a within-model sign test rather than
a threshold count.

\paragraph{Nothing here is significant.} Eight of eleven above median is a
one-sided binomial $p = 0.113$; the best permutation $p$ across all rankers is
$0.016$ and survives no correction for the fourteen comparisons made. Restricting to
the four trustworthy coordinates raises hydrology's AUC to $0.874$ ($p=0.003$)
but at $n=4$ that is an observation, not a result.

\subsection{What eleven points can and cannot support}

They can reject: a ranker scoring at or below chance on these locations is not
usable in Al~Madinah, which excludes Naive Bayes ($0.431$), logistic regression
($0.478$), SVM-RBF ($0.488$) and the decision tree ($0.515$). They cannot select: the spread among the four boosting variants
($0.635$--$0.644$) is far narrower than what best-of-fourteen selection produces under
a null at this sample size. And they cannot estimate precision at all, because
an unreported street is not a dry street --- the same presence-only structure
analysed in Section~\ref{sec:pu}, now with three orders of magnitude less data.

The honest summary is that the transfer is encouraging and unvalidated. The
measurement that would settle it is a single rain event observed jointly with
the municipality, which converts eleven retrospective points into a labelled
sample of the size this evaluation actually requires.

\section{Discussion}\label{sec:discussion}

\textbf{The field optimises the wrong axis.} Our results reproduce, under a
harder protocol, what \citet{zhang2025} found for sampling: the spread across
\num{14} model families is $0.099$ skill, while a single feature-representation
change worth $0.0137$ exceeds the gap between the best and third-best model, and
per-city normalisation of terrain is worth more than replacing the learner
entirely. A literature that compares architectures on balanced samples is
measuring the axis with the least available variance.

\textbf{Absolute terrain features are the wrong representation for transfer.}
That terrain scores below random across cities while working within them is, to
our knowledge, unreported, and it has a clean mechanistic explanation and a
cheap fix. It also reframes \citet{agonafir2022ml}, who found building factors
outweighing elevation in New York: we find the same ordering, and show it
strengthens rather than weakens when the evaluation is made harder.

\textbf{Our scores are not comparable to the literature's, and that is the
point.} Table~\ref{tab:lit} reports AUC up to $0.985$; we report ROC-AUC
$0.649$. \citet{agostini2024}, modelling the same New York feed at natural
prevalence with explicit under-reporting, report $0.642$. Two independent
analyses of the real problem agree; the balanced-sample results are measuring
something else.

\textbf{Implication for Al Madinah.} The honest transfer expectation is not the
mean across cities but the value for the city that most resembles the target.
Dallas gives $\skl = +0.027$. A deployment claim should be stated at that level,
with per-segment coverage flags, and as a \emph{rank} rather than a probability
--- a model transferred across domains is not calibrated, and an unlabelled city
offers nothing to calibrate against.

% ══════════════════════════════════════════════════════════ LIMITS
\section{Limitations}\label{sec:limits}

\begin{enumerate}\itemsep2pt
  \item \textbf{Four cities is still small.} Our own Table~\ref{tab:flip} shows
    rankings inverting between $n{=}2$ and $n{=}4$; we cannot exclude further
    change at $n{=}6$. The top four models sit within $0.018$ of each other.
  \item \textbf{All four cities are in the United States.} Transfer to a Saudi
    city crosses climate, drainage standard and reporting culture simultaneously,
    and no held-out city in this benchmark tests that.
  \item \textbf{Labels are complaints, not observations.} \S\ref{sec:pu}
    quantifies the resulting incompleteness but cannot remove it, and the
    San Francisco estimate shows the correction's own assumption failing.
  \item \textbf{No drainage-infrastructure layer.} \S\ref{sec:pu} shows that
    terrain ponding and drainage failure are inseparable with these features.
  \item \textbf{$\SI{30}{\metre}$ elevation inside a dense city} resolves
    catchment topography, not kerb-level camber; the DEM is a surface model
    including buildings.
  \item \textbf{Al Madinah is evaluated, but not validated.} The distributional
    numbers in \S\ref{sec:ood} describe feature space, not accuracy. The eleven
    locations in \S\ref{sec:target} are a genuine ground-truth test, but an
    underpowered one: they are positives only, assembled retrospectively from
    heterogeneous sources, and four of the eleven carry coordinates we would
    defend without qualification. They can exclude a ranker; they cannot certify
    one, and no accuracy figure for Al Madinah should be quoted from them.
\end{enumerate}

% ══════════════════════════════════════════════════════════ CONCLUSION
\section{Conclusion}\label{sec:conclusion}

We benchmarked \num{14} model families for street-flood susceptibility across
four cities at natural prevalence under spatial blocking, with cross-city
transfer as the declared endpoint. Gradient boosting leads, but the margin is
small, the ranking is unstable at the numbers of held-out cities this field can
realistically assemble, and the largest single improvement we found came from
changing how terrain features are represented rather than which model consumes
them. Absolute terrain descriptors do not transfer between cities and can be
repaired by per-city percentile normalisation, which needs no labels and is
therefore available in exactly the unlabelled setting that motivates transfer.
Presence-only corrections quantify a real and substantial reporting gap but
cannot, by construction, improve a ranking. And distributional distance did not
predict transfer quality in our hands.

Carried to the target city and tested on eleven hand-assembled flood locations,
the whole apparatus of transfer proved indistinguishable from four lines of
hydrology with nothing fitted: AUC $0.685$ against $0.674$, with the transferred
network ahead on lift. We report that as a tie rather than a defeat for
transfer, because at $n{=}11$ nothing else is defensible. But a tie is decisive
in practice. One side needs \num{131298} foreign labels, a training pipeline and
a frozen preprocessing contract; the other needs four rasters and a sign
convention, and can be audited by the engineer who has to act on it. The
burden of proof belongs on the complicated method, and here it was not
discharged.

For a city with no flood records, the defensible protocol is: never train on it;
report leave-one-city-out rather than within-city cross-validation; quote the
held-out city that most resembles the target rather than the mean; normalise
terrain within the target city; publish a coverage flag alongside every score;
output a rank, not a probability; and benchmark against an untrained physical
index before accepting a transferred model, because it may not be earning its
complexity.

\paragraph{Reproducibility.} Feature construction, the model registry, all
evaluation protocols and every results file are released with the manuscript.
Library versions are pinned: an XGBoost minor-version change alone moved PR-AUC
from $0.340$ to $0.342$ in our own pipeline, which is small but larger than
several of the model gaps in Table~\ref{tab:board}.

% ══════════════════════════════════════════════════════════ BACK MATTER
\section*{Data and code availability}
All inputs are public. The flood complaints come from the four cities' open-data
portals (NYC 311, Chicago 311, SF~311 \texttt{vw6y-z8j6}, Dallas
\texttt{d7e7-envw}); elevation from Copernicus DEM GLO-30 (ESA, free licence);
street geometry from OpenStreetMap (ODbL); precipitation from the Open-Meteo
archive and NASA POWER. The eleven Al~Madinah flood locations were assembled by
the author from official statements, press reports and resident accounts; the
register, including each coordinate's provenance and a quality grade, is released
with the code so that the four \emph{trustworthy} points can be separated from
the seven weaker ones by any reader.

Feature construction, the model registry, all fourteen fitted models with their
preprocessing contracts, every evaluation protocol and every results file are
released. Library versions are pinned: an XGBoost minor-version change alone
moved PR-AUC from $0.340$ to $0.342$ in our own pipeline.

\section*{Declaration of competing interests}
The author declares no competing financial or personal interests.

\section*{Funding}
This work received no external funding. It was developed for Barmajan
Al~Munawwarah 2026, which provided no research funding and had no role in study
design, analysis, or the decision to publish.

\section*{Use of the target city's data}
No personal data was processed. The Al~Madinah incident register contains street
locations and dates drawn from published sources; individual reporters are not
identified, and resident accounts are recorded at street granularity without
attribution.

\bibliographystyle{plainnat}
\bibliography{refs}

\end{document}
